% META: {"id": "TSPE-DERCONV-COURS-07-TC", "chapitre": "TSPE-DERIVATION-CONVEXITE", "type_objet": "cours", "sous_type": "td_contextualise", "titre": "TD contextualise -- Optimisation d'une boite de conserve", "temps_gabarit": 7, "capacites": ["C1", "C2", "C3", "C6"], "mode_creation": "ex_nihilo", "sources_inspiration": [], "status": "generated"}

% BEGIN-VERIFY
% from sympy import *
% x, r, h = symbols('x r h', positive=True)
% # Volume V = pi*r^2*h = 500 cm^3, h = 500/(pi*r^2)
% # Surface S(r) = 2*pi*r^2 + 2*pi*r*h = 2*pi*r^2 + 1000/r
% S = 2*pi*x**2 + 1000/x
% Sp = diff(S, x)
% assert simplify(Sp - (4*pi*x - 1000/x**2)) == 0
% # S'(x) = 0 => 4*pi*x^3 = 1000 => x^3 = 250/pi
% r0 = (250/pi)**Rational(1,3)
% assert simplify(Sp.subs(x, r0)) == 0
% # S''(x) = 4*pi + 2000/x^3
% Spp = diff(S, x, 2)
% assert simplify(Spp - (4*pi + 2000/x**3)) == 0
% # S''(r0) > 0 (convexe)
% assert Spp.subs(x, r0) > 0
% # r0 ~ 4.30 cm
% import math
% r0_num = (250/math.pi)**(1/3)
% assert abs(r0_num - 4.30) < 0.01
% h0_num = 500/(math.pi * r0_num**2)
% assert abs(h0_num - 2*r0_num) < 0.01
% END-VERIFY

\section*{Temps 7 -- TD contextualise : Optimisation d'une boite de conserve}

\textbf{Contexte.} Un fabricant souhaite produire une boite de conserve cylindrique de volume $V = 500$~cm$^3$. Il cherche a minimiser la quantite de metal utilisee, c'est-a-dire la surface totale du cylindre. On note $r$ le rayon de la base (en cm) et $h$ la hauteur.

\bigskip

%---------------------------------------------------------------
% PARTIE A — Mise en equation (C1, C2)
%---------------------------------------------------------------

\begin{exercice}{TSPE-DERCONV-COURS-07-TC-EX1}{2}{15}

\textbf{Partie A -- Mise en equation}

\begin{enumerate}
  \item Exprimer la contrainte de volume : $\pi r^2 h = 500$. En deduire $h$ en fonction de $r$.
  \item Montrer que la surface totale $S(r) = 2\pi r^2 + 2\pi r h$ s'ecrit $S(r) = 2\pi r^2 + \dfrac{1000}{r}$ pour $r > 0$.
  \item Calculer $S'(r)$.
  \item Resoudre $S'(r) = 0$ et montrer que le rayon optimal est $r_0 = \left(\dfrac{250}{\pi}\right)^{1/3} \approx 4{,}30$~cm.
\end{enumerate}

\end{exercice}

\begin{corrige}{TSPE-DERCONV-COURS-07-TC-EX1}

\textbf{Question 1.} $h = \dfrac{500}{\pi r^2}$.

\textbf{Question 2.} $S(r) = 2\pi r^2 + 2\pi r \cdot \dfrac{500}{\pi r^2} = 2\pi r^2 + \dfrac{1000}{r}$.

\commentaireMarge{On utilise la formule de derivation de $u^{-1}$ : $(1/r)' = -1/r^2$.}

\textbf{Question 3.} $S'(r) = 4\pi r - \dfrac{1000}{r^2}$.

\textbf{Question 4.} $S'(r) = 0 \iff 4\pi r = \dfrac{1000}{r^2} \iff 4\pi r^3 = 1000 \iff r^3 = \dfrac{250}{\pi}$, soit $r_0 = \left(\dfrac{250}{\pi}\right)^{1/3} \approx 4{,}30$~cm.

\end{corrige}

\bigskip

%---------------------------------------------------------------
% PARTIE B — Convexite et minimum (C3, C6)
%---------------------------------------------------------------

\begin{exercice}{TSPE-DERCONV-COURS-07-TC-EX2}{2}{15}

\textbf{Partie B -- Convexite et minimum global}

\begin{enumerate}
  \item Calculer $S''(r)$ et montrer que $S''(r) > 0$ pour tout $r > 0$.
  \item En deduire que $S$ est convexe sur $]0, +\infty[$.
  \item Justifier que $r_0$ realise bien un minimum global de $S$.
  \item Calculer la hauteur correspondante $h_0$ et verifier que $h_0 = 2r_0$ (le diametre est egal a la hauteur).
\end{enumerate}

\end{exercice}

\begin{corrige}{TSPE-DERCONV-COURS-07-TC-EX2}

\textbf{Question 1.} $S''(r) = 4\pi + \dfrac{2000}{r^3}$. Pour $r > 0$, on a $\dfrac{2000}{r^3} > 0$ et $4\pi > 0$, donc $S''(r) > 0$.

\textbf{Question 2.} Puisque $S''(r) > 0$ pour tout $r > 0$, la fonction $S$ est convexe sur $]0, +\infty[$.

\textbf{Question 3.} Pour une fonction convexe, tout minimum local est un minimum global. Or $S'(r_0) = 0$ et $S''(r_0) > 0$, donc $r_0$ est un minimum local, et par convexite, c'est un minimum global.

\textbf{Question 4.} $h_0 = \dfrac{500}{\pi r_0^2} = \dfrac{500}{\pi \cdot (250/\pi)^{2/3}} = \dfrac{500}{\pi^{1/3} \cdot 250^{2/3}}$.

Or $2r_0 = 2 \cdot (250/\pi)^{1/3}$. Verifions : $2r_0 = \dfrac{2 \cdot 250^{1/3}}{\pi^{1/3}}$ et $h_0 = \dfrac{500}{\pi \cdot 250^{2/3}/\pi^{2/3}} = \dfrac{500 \cdot \pi^{2/3}}{\pi \cdot 250^{2/3}} = \dfrac{500}{\pi^{1/3} \cdot 250^{2/3}} = \dfrac{2 \cdot 250}{\pi^{1/3} \cdot 250^{2/3}} = \dfrac{2 \cdot 250^{1/3}}{\pi^{1/3}} = 2r_0$.

Donc $h_0 = 2r_0 \approx 8{,}60$~cm.

\end{corrige}
